\subsection{Conformal Prediction}
% \noindent - framework\\
% - coverage guarantee\\
% - exchnagebility

\subsubsection{Setup and nonconformity scores}
\noindent - define exchangeability, define s(x,y) -> explain the intuition

\subsubsection{prediction set and coverage guarantee} \noindent - define C(x) 
%state Theorem 2.1 formally with both bounds

\subsubsection{Split conformal prediction} 
\noindent - why we use it, the train/calibration/test separation, validity still holds

% \subsubsection{Marginal vs conditional coverage}
% \noindent - the distinction, why it matters for imbalanced problems, formal statement of both guarantees
% Marginal guarantee: holds on average over all test points
% Why this can hide severe under-coverage of minority classes — numerical example
% Formal statement of class-conditional coverage: P(Y∈C(X)∣Y=y)≥1−αP(Y \in C(X) \mid Y = y) \geq 1 - \alpha
% P(Y∈C(X)∣Y=y)≥1−α for all yy
% y
% This motivates the class-conditional t^\hat{t}
% t^ developed in section 3.4.5
